\documentclass[11pt]{article}
\usepackage[margin=1in]{geometry}
\usepackage{booktabs}
\usepackage{amsmath}
\usepackage{siunitx}
\usepackage{hyperref}
\usepackage{caption}

\title{Weighting Scheme: Inverse Average-Correlation}
\author{MetaFVG Diversified Mean-Reversion Portfolio}
\date{}

\begin{document}
\maketitle

\section{Scheme Description}

The 14-instrument MetaFVG mean-reversion portfolio (Spearman-gated, \texttt{fixed} sizing
variant) is deliberately assembled to be fundamentally diversified across central
banks and currencies: AUDUSD, EURUSD, GBPUSD, NZDUSD, USDCAD, USDCHF, USDJPY,
US500, GER40, JP225, HK50, USDSGD, EURNOK, USDZAR. It is currently blended with
equal weights ($1/14$ each). This report tests \textbf{Inverse Average-Correlation}
weighting (\texttt{inverse\_avg\_corr}): the most literal reading of ``target
uncorrelated results'' --- weight each instrument inversely proportional to how
correlated it is, on average, with the rest of the book.

For each instrument $s$, let $\rho_{s,x}$ denote the pairwise Pearson correlation
of daily P\&L returns between $s$ and instrument $x$ (computed pairwise-complete
over the full return history). The average correlation of $s$ with the rest of
the book is

\[
\overline{\rho}_s = \frac{1}{n-1} \sum_{x \neq s} \rho_{s,x}, \qquad n = 14.
\]

Because the portfolio was already fundamentally diversified by design, these
average correlations are small and can be near zero or slightly negative ---
naively inverting such a value blows up. A floor is applied before inversion:

\[
d_s = \max(\overline{\rho}_s,\ 0.02).
\]

Raw weights are then

\[
\tilde{w}_s \propto \frac{1}{d_s}, \qquad
w_s^{(0)} = \frac{\tilde{w}_s}{\sum_x \tilde{w}_x}.
\]

Finally, a concentration cap is applied at $3\times$ the equal weight
($3/14 \approx 0.2143$): any weight exceeding the cap is clipped to it, and the
remaining (uncapped) weights are renormalized proportionally so that all weights
still sum to 1.

\textbf{Safeguards triggered:} the floor bound for \textbf{9 of 14} instruments ---
EURNOK, EURUSD, HK50, JP225, USDCAD, USDCHF, USDJPY, USDSGD, USDZAR --- all of
which had raw average correlation below 0.02 (three of them, EURNOK, USDCHF and
HK50, were in fact very close to zero or negative). These instruments were all
floored to the same denominator (0.02) and therefore received an identical final
weight ($\approx 0.0811$, or $1.136\times$ equal weight). The 3x concentration
cap ($0.2143$) did \textbf{not} bind for any instrument -- the highest resulting
weight was only $\approx 1.14\times$ equal weight, well under the cap.

\section{Weights}

\begin{table}[h]
\centering
\caption{Per-instrument average correlation, resulting weight, and weight relative
to equal weight ($1/14 \approx 0.0714$), sorted descending by weight. F = floor
bound, C = cap bound.}
\begin{tabular}{lrrrc}
\toprule
Symbol & Avg.\ Corr.\ ($\overline{\rho}_s$) & Weight & Weight / (1/14) & Safeguard \\
\midrule
EURUSD & 0.0196 & 0.0811 & 1.14$\times$ & F \\
USDCAD & 0.0168 & 0.0811 & 1.14$\times$ & F \\
USDCHF & 0.0004 & 0.0811 & 1.14$\times$ & F \\
USDJPY & 0.0056 & 0.0811 & 1.14$\times$ & F \\
JP225  & 0.0188 & 0.0811 & 1.14$\times$ & F \\
HK50   & 0.0016 & 0.0811 & 1.14$\times$ & F \\
USDSGD & 0.0120 & 0.0811 & 1.14$\times$ & F \\
EURNOK & $-$0.0153 & 0.0811 & 1.14$\times$ & F \\
USDZAR & 0.0070 & 0.0811 & 1.14$\times$ & F \\
NZDUSD & 0.0265 & 0.0613 & 0.86$\times$ & -- \\
US500  & 0.0276 & 0.0588 & 0.82$\times$ & -- \\
AUDUSD & 0.0307 & 0.0528 & 0.74$\times$ & -- \\
GER40  & 0.0313 & 0.0519 & 0.73$\times$ & -- \\
GBPUSD & 0.0360 & 0.0450 & 0.63$\times$ & -- \\
\bottomrule
\end{tabular}
\end{table}

\section{Correlation Impact}

Weighted average pairwise correlation is
$\left(\sum_{i \neq j} w_i w_j \rho_{ij}\right) / \left(\sum_{i \neq j} w_i w_j\right)$;
for equal weights this reduces to the simple average of the off-diagonal
correlation matrix ($\approx 0.0156$). Diversification ratio is
$\mathrm{DR} = \left(\sum_i w_i \sigma_i\right) / \sigma_{\text{portfolio}}$,
using each instrument's own-sample standard deviation of daily returns and the
standard deviation of the final blended daily-return series.

\begin{table}[h]
\centering
\caption{Correlation impact: equal weight vs.\ inverse-average-correlation weight.}
\begin{tabular}{lrr}
\toprule
Metric & Equal Weight & Inverse Avg.\ Corr. \\
\midrule
Weighted avg.\ pairwise correlation & 0.0156 & 0.0118 \\
Diversification ratio (DR) & 2.081 & 2.110 \\
\bottomrule
\end{tabular}
\end{table}

\section{Results}

Both blends use per-day NaN-aware renormalization: on days where some instruments
have no closed trade (NaN return), the scheme's weights are renormalized to sum to
1 over only the instruments with a non-NaN return that day. Days with all-NaN
returns are dropped. Both series contain 1{,}661 daily observations after this
filtering. Quantstats metrics are computed on the resulting equity curve (daily
resampled).

\begin{table}[h]
\centering
\caption{Performance metrics: equal weight vs.\ inverse-average-correlation weight.}
\begin{tabular}{lrr}
\toprule
Metric & Equal Weight & Inverse Avg.\ Corr. \\
\midrule
Sharpe & 0.518 & 0.483 \\
Sortino & 0.817 & 0.760 \\
Calmar & 0.343 & 0.315 \\
CAGR & 0.0183\% & 0.0168\% \\
Max Drawdown & $-5.33\%$ & $-5.33\%$ \\
Volatility (ann., \%) & 3.528\% & 3.473\% \\
Tail Ratio & 1.239 & 1.288 \\
Ulcer Index & 0.000192 & 0.000173 \\
Kelly & 0.0668 & 0.0610 \\
VaR (\%) & $-0.358\%$ & $-0.353\%$ \\
CVaR (\%) & $-0.715\%$ & $-0.692\%$ \\
Skew & 1.545 & 1.544 \\
Recovery Factor & 2.260 & 2.075 \\
\midrule
Win Rate & 30.89\% & 30.52\% \\
Best Day & 0.0291\% & 0.0291\% \\
Worst Day & $-0.0239\%$ & $-0.0239\%$ \\
Total Return & 0.1204\% & 0.1105\% \\
\bottomrule
\end{tabular}
\end{table}

\section{Takeaway}

Inverse average-correlation weighting achieved the intended correlation effect:
weighted average pairwise correlation fell from 0.0156 (equal weight) to 0.0118,
and the diversification ratio improved modestly from 2.081 to 2.110. Neither
safeguard was actually a problem in practice --- the floor caught 9 of 14
instruments (clustering them at an identical weight), while the 3x cap never
bound, meaning the floor did all the real work of shaping this allocation.
However, the reduction in realized correlation did \emph{not} translate into
better risk-adjusted returns: Sharpe fell from 0.518 to 0.483 and Sortino,
Calmar, Kelly, and Recovery Factor all declined in tandem. The reason is visible
in the individual-instrument statistics: the scheme's three lowest final weights
went to GBPUSD, AUDUSD and NZDUSD (annualized single-instrument Sharpe 0.76,
0.66 and 0.76 respectively --- three of the book's best performers), while the
floor-driven cluster of nine equally-weighted instruments includes EURUSD,
USDCAD, USDSGD, EURNOK and USDZAR, all of which have \emph{negative}
single-instrument Sharpe on their own. Low average correlation with the rest of
the book and being a strong individual performer are simply different
properties, and this scheme optimizes only for the former. On this dataset it
trades away Sharpe for a small diversification-ratio gain --- a bad trade for a
book that was already reasonably diversified by construction.

\end{document}
